% Chapitre 2 — séance 1 : binaire, masque, ET logique, taille de bloc
% Accroche : « pourquoi les liaisons routeur étaient-elles en /30 ? »
% Compilation : make 02 (versions élève et prof) — voir README.md
\documentclass[aspectratio=169,10pt]{beamer}
\usepackage{ciel-beamer}

\title{Pourquoi /30 ?}
\subtitle{BTS CIEL IR · 2\textsuperscript{e} année · Chapitre 2 · Séance 1}
\author{Réseau — Subnetting et VLSM}
\date{Binaire, masque, ET logique, taille de bloc}
\chapitre{Ch. 2 · Subnetting}

\begin{document}

% ------------------------------------------------------------------ 1
\begin{frame}[plain]
  \titlepage
\end{frame}

% ------------------------------------------------------------------ 2 accroche
\begin{frame}
\frametitle{Pourquoi des /30 entre les routeurs ?}
\framesubtitle{Accroche · le réseau fil rouge}
\begin{columns}[T]
\column{.42\linewidth}
\begin{carte}[tag=Topologie]{}
\centering
\begin{tikzpicture}[x=1cm,y=1cm]
  \node[parefeu] (FW) at (0,0) {};
  \node[routeur] (R2) at (1.75,0) {};
  \node[routeur] (R1) at (3.5,0) {};
  \node[switch]  (S1) at (3.5,-1.25) {};
  \draw[lien] (FW) -- node[below=1pt,adr]{.228/30} (R2);
  \draw[lien] (R2) -- node[below=1pt,adr]{.224/30} (R1);
  \draw[trunk] (R1) -- (S1);
  \node[equip,above=1pt of FW]{FW};
  \node[equip,above=1pt of R2]{R2};
  \node[equip,above=1pt of R1]{R1};
  \node[equip,left=2pt of S1,align=right]{VLAN 10, 20, 30};
\end{tikzpicture}
\end{carte}
\column{.53\linewidth}
\begin{carte}[tag=Plan 192.168.10.0/24]{}
\footnotesize
\setlength{\tabcolsep}{4pt}\begin{tabular}{@{}l M M@{}}
Production   & 192.168.10.0/25   & .126\\
Admin.       & 192.168.10.128/26 & .190\\
Wi-Fi invité & 192.168.10.192/27 & .222\\
R1–R2        & 192.168.10.224/30 & --\\
R2–FW        & 192.168.10.228/30 & --\\
\end{tabular}
\end{carte}
\end{columns}
\vfill
\begin{callout}[or]
\cielfontlabel Pourquoi /30 et pas /24 ou /27 ? En fin de séance, vous saurez le démontrer.
\end{callout}
\end{frame}

% ------------------------------------------------------------------ 3 32 bits
\begin{frame}
\frametitle{Une adresse IPv4, c'est 32 bits}
\framesubtitle{Rappel · 4 octets de 8 bits}
\begin{carte}[tag=Exemple · adresse de R2 côté R1]{}
\centering
{\cielfontlabel\Large 192\,·\,168\,·\,10\,·\,\textcolor{or!70!black}{226}}\par\medskip
\poids\par
\octet{1,1,1,0,0,0,1,0}\par\medskip
\small 226 = 128 + 64 + 32 + 2
\end{carte}
\vfill
\begin{columns}[T]
\column{.48\linewidth}
\begin{callout}[rose]\small Un octet vaut de 0 (00000000) à 255 (11111111).\end{callout}
\column{.48\linewidth}
\begin{callout}[vert]\small Chaque bit a un poids : 128, 64, 32, 16, 8, 4, 2, 1.\end{callout}
\end{columns}
\end{frame}

% ------------------------------------------------------------------ 4 conversion
\begin{frame}
\frametitle{Décimal vers binaire : on retire les poids}
\framesubtitle{Méthode}
\begin{columns}[T]
\column{.42\linewidth}
\begin{carte}[tag=Méthode]{}
\small
\begin{enumerate}
\item Partir de 128.
\item Le poids tient dans le reste ? \\ bit = 1 et on le retire.
\item Sinon bit = 0.
\item Poids suivant, jusqu'à 1.
\end{enumerate}
\end{carte}
\column{.55\linewidth}
\begin{carte}[tag=À vous,rose]{}
\small
\begin{tabular}{@{}l l@{}}
190 & \reponse[30mm]{10111110 \;(128+32+16+8+4+2)}\\[2mm]
252 & \reponse[30mm]{11111100 \;(128+64+32+16+8+4)}\\[2mm]
224 & \reponse[30mm]{11100000 \;(128+64+32)}\\[2mm]
128 & \reponse[30mm]{10000000}\\
\end{tabular}
\end{carte}
\end{columns}
\end{frame}

% ------------------------------------------------------------------ 5 masque
\begin{frame}
\frametitle{Le masque : des 1, puis des 0}
\framesubtitle{Masque · notation /n}
\begin{carte}[tag=/30]{30 bits à 1 = partie réseau, 2 bits à 0 = partie hôte}
\centering\small
\begin{tabular}{@{}c c c c@{}}
\cielfontlabel 11111111 & \cielfontlabel 11111111 & \cielfontlabel 11111111 & \octet[6]{1,1,1,1,1,1,0,0}\\[1mm]
\cielfontlabel 255 & \cielfontlabel 255 & \cielfontlabel 255 & \cielfontlabel 252\\
\end{tabular}\par\medskip
\ip{/30} \;=\; \ip{255.255.255.252}
\end{carte}
\vfill
\begin{callout}[or]\small
Les 1 sont toujours \textbf{à gauche} et \textbf{contigus} : 255.255.255.\textbf{250} n'est pas un masque.
\end{callout}
\end{frame}

% ------------------------------------------------------------------ 6 ET logique
\begin{frame}
\frametitle{Adresse ET masque = adresse du réseau}
\framesubtitle{ET logique · bit à bit}
\begin{columns}[T]
\column{.28\linewidth}
\begin{carte}[tag=Table]{}
\centering\cielfontlabel
\begin{tabular}{@{}ccc@{}}
a & b & a ET b\\\midrule
0 & 0 & 0\\ 0 & 1 & 0\\ 1 & 0 & 0\\ 1 & 1 & 1\\
\end{tabular}
\end{carte}
\column{.69\linewidth}
\begin{carte}[tag=R2 · dernier octet]{}
\small
\begin{tabular}{@{}l l@{}}
226 & \octet{1,1,1,0,0,0,1,0}\\[1mm]
252 & \octet[6]{1,1,1,1,1,1,0,0}\\[1mm]\midrule
ET  & \reponse[44mm]{\octet[6]{1,1,1,0,0,0,0,0}\; = 224}\\
\end{tabular}\par\medskip
Réseau de R2 : \reponse[30mm]{\ip{192.168.10.224/30}}
\end{carte}
\end{columns}
\vfill
\begin{callout}[vert]\small Un 0 dans le masque efface le bit : il ne reste que la partie réseau.\end{callout}
\end{frame}

% ------------------------------------------------------------------ 7 taille de bloc
\begin{frame}
\frametitle{Taille de bloc = 256 − octet du masque}
\framesubtitle{Le raccourci}
\begin{carte}[tag=/30]{256 − 252 = 4 : les réseaux avancent de 4 en 4}
\centering
\begin{tikzpicture}[x=1cm]
  \foreach \v [count=\i] in {.224,.228,.232,.236,.240}{
    \node[fill=cielicon,text=cieliconfg,rounded corners=1mm,minimum width=17mm,
      minimum height=7mm,font=\cielfontlabel\small] at (\i*2.1,0) {\v};}
  \foreach \i in {1,...,4}
    \draw[-{Triangle[length=2pt]},cielicon,line width=.8pt] (\i*2.1+.9,0) -- (\i*2.1+1.2,0);
\end{tikzpicture}\par\medskip
\small Un réseau commence toujours sur un multiple de la taille de bloc.
\end{carte}
\vfill
\begin{columns}[T]
\column{.48\linewidth}
\begin{callout}[or]\small /27 : 256 − 224 = \reponse{32} \\ réseaux .0, .32, .64 … .192, .224\end{callout}
\column{.48\linewidth}
\begin{callout}[rose]\small /26 : 256 − 192 = \reponse{64} \\ réseaux .0, .64, .128, .192\end{callout}
\end{columns}
\end{frame}

% ------------------------------------------------------------------ 8 réponse
\begin{frame}
\frametitle{Dans un /30 : 4 adresses, 2 utilisables}
\framesubtitle{Réponse à l'accroche}
\begin{carte}[tag=192.168.10.224/30]{}
\centering
\begin{tikzpicture}[x=1cm]
  \foreach \v/\r/\s [count=\i] in {.224/réseau/0,.225/{R1 G0/1}/1,.226/{R2 G0/0}/1,.227/diffusion/0}{
    \ifnum\s=1
      \node[fill=cielicon,text=cieliconfg,rounded corners=1mm,minimum width=24mm,
        minimum height=9mm,font=\cielfontlabel\normalsize] (c\i) at (\i*2.7,0) {\v};
    \else
      \node[draw=cielicon,line width=.8pt,text=cielicon,rounded corners=1mm,minimum width=24mm,
        minimum height=9mm,font=\cielfontlabel\normalsize] (c\i) at (\i*2.7,0) {\v};
    \fi
    \node[font=\cielfontlabel\scriptsize,text=cielicon,below=1mm of c\i] {\r};}
\end{tikzpicture}
\end{carte}
\vfill
\begin{columns}[T]
\column{.48\linewidth}
\begin{callout}[vert]\small Une liaison entre deux routeurs = exactement 2 interfaces. \\ Le /30 n'en gaspille aucune.\end{callout}
\column{.48\linewidth}
\begin{callout}[rose]\small Avec un /24 pour la même liaison : \reponse{252} adresses perdues.\end{callout}
\end{columns}
\end{frame}

% ------------------------------------------------------------------ 9 exercice
\begin{frame}
\frametitle{À vous : réseau et diffusion}
\framesubtitle{Entraînement · 5 minutes}
\begin{carte}[tag=Exercice,rose]{}
\small
\renewcommand{\arraystretch}{1.5}
\begin{tabular}{@{}M l l l@{}}
\textrm{\cielfontlabel Adresse} & \cielfontlabel Bloc & \cielfontlabel Réseau & \cielfontlabel Diffusion\\\midrule
192.168.10.229/30 & \reponse[14mm]{4}  & \reponse[31mm]{\ip{192.168.10.228}} & \reponse[31mm]{\ip{192.168.10.231}}\\
192.168.10.200/27 & \reponse[14mm]{32} & \reponse[31mm]{\ip{192.168.10.192}} & \reponse[31mm]{\ip{192.168.10.223}}\\
192.168.10.150/26 & \reponse[14mm]{64} & \reponse[31mm]{\ip{192.168.10.128}} & \reponse[31mm]{\ip{192.168.10.191}}\\
\end{tabular}
\end{carte}
\vfill
\begin{callout}[or]\small Diffusion = adresse du réseau suivant − 1.\end{callout}
\end{frame}

% ------------------------------------------------------------------ 10 finale inversée
\begin{finale}
\begin{frame}
\frametitle{À retenir}
\framesubtitle{Séance 1}
\begin{columns}[T]
\column{.3\linewidth}
\begin{carte}[equal height group=finale,tag=01]{8 bits}
\small Poids 128, 64, 32, 16, 8, 4, 2, 1. Un octet va de 0 à 255.
\end{carte}
\column{.3\linewidth}
\begin{carte}[equal height group=finale,tag=02]{IP ET masque}
\small Le ET bit à bit donne l'adresse du réseau.
\end{carte}
\column{.3\linewidth}
\begin{carte}[equal height group=finale,tag=03]{256 − octet}
\small Taille de bloc : les réseaux en sont des multiples. /30 = 4 adresses.
\end{carte}
\end{columns}
\vfill
\begin{callout}[or]\small \cielfontlabel Prochaine séance : découper 192.168.10.0/24 en sous-réseaux égaux.\end{callout}
\end{frame}
\end{finale}

\end{document}
